\documentclass[twocolumn]{article}
\usepackage[utf8]{inputenc}
\usepackage[round]{natbib}
\begin{document}

\section{Introduction}
Knowledge editing changes what a language model believes about the world, and evaluating it requires benchmarks whose metrics are defined precisely enough that two papers reporting the same number mean the same thing.
This paragraph is padding so that the two columns fill and a sentence has to cross from the first column into the second one, which is the case the scanner must survive without splitting the sentence in two.
We repeat the point once more to push the text further down the first column of the first page of this short document.

\citet{zhong2023mquake} define multi-hop accuracy as correct if any of three paraphrased questions is answered.
MQuAKE~\citep{zhong2023mquake} evaluates whether an edit propagates to its multi-hop consequences.
Temporal benchmarks \citep{jang2022temporalwiki,zhong2023mquake} track how facts change over time.
Mutual information follows the definition in \citet{cover2006elements}.
The scaling result of \citet{smith2020first} was later revisited under distribution shift \citep{smith2020second}.
We follow the evaluation protocol of \citet{bao2023tallrec} for all ranking runs.
A claim attributed to an unknown source (Nobody et al., 2019) has no entry in the list.
\bibliographystyle{plainnat}
\bibliography{refs}
\appendix
\section{Extra Results}
The appendix repeats the protocol of \citet{bao2023tallrec} on a second dataset.
\end{document}
